\answer
\begin{enumerate}

%%--- 章末問題 6.1 ---
\item[\ref{ex6.1}]
式\ref{eq6.4}より，電圧が$+10$\,Vの間，$i_m$は一定の傾き
\begin{equation*}
\frac{di_m}{dt} = \frac{v_1}{L_m} = \frac{10}{1\times10^{-3}}
= 1.0\times10^{4}\,\mathrm{A/s}
\end{equation*}
で増える。半周期$1.0$\,$\mu$sでの増加量は
\begin{equation*}
\mathit{\Delta}I_m = \frac{v_1}{L_m}\,\mathit{\Delta}t
= 1.0\times10^{4} \times 1.0\times10^{-6} = 10\,\mathrm{mA}
\end{equation*}
電圧が$-10$\,Vの半周期には同じ傾きで減る。
方形波は正負対称でボルト秒平衡が自動的に満たされるから，
$i_m$は$-5$\,mAから$+5$\,mAの間を往復する三角波になる。
電圧が方形波，電流が三角波 ― 5章のインダクタ電流と同じ形である。

%%--- 章末問題 6.2 ---
\item[\ref{ex6.2}]
(a) 式\ref{eq6.12}より
\begin{equation*}
\rev{V_{\mathrm{out}} = \frac{N_2}{N_1}DV_{\mathrm{in}} = \frac{1}{6}\times0.375\times48 = 3.0\,\mathrm{V}}
\end{equation*}
(b) \rev{周期は$T = 1/f = 10\,\mathrm{\mu s}$。
オフ期間$(1-D)T = 6.25\,\mathrm{\mu s}$の間，$v_L = -V_{\mathrm{out}}$だから}
\begin{equation*}
\rev{\mathit{\Delta}I_L = \frac{V_{\mathrm{out}}}{L}(1-D)T \le 0.2\,\mathrm{A}}
\end{equation*}
これを$L$について解くと
\begin{equation*}
\rev{L \ge \frac{V_{\mathrm{out}}(1-D)T}{\mathit{\Delta}I_L}
= \frac{3.0\times6.25\times10^{-6}}{0.2} \approx 94\,\mathrm{\mu H}}
\end{equation*}
となる。たとえば$100\,\mathrm{\mu H}$を選べばよい。
リプルの式は5章の降圧コンバータと同形であり，
出力段が降圧コンバータそのものであることの表れである。

%%--- 章末問題 6.3 ---
\item[\ref{ex6.3}]
(a) 式\ref{eq6.15}に$N_3/N_1 = 1/2$を代入して
\begin{equation*}
D \le \frac{1}{1 + 1/2} = \frac{2}{3}
\end{equation*}
(b) リセット期間のスイッチ電圧は電源電圧と1次巻線電圧の和だから
\begin{equation*}
\rev{v_{sw} = V_{\mathrm{in}} + \frac{N_1}{N_3}V_{\mathrm{in}} = V_{\mathrm{in}} + 2V_{\mathrm{in}} = 3V_{\mathrm{in}}}
\end{equation*}
(c) 利点：デューティ比の上限が$0.5$から$2/3$へ広がり，
入力電圧が低いときも必要な出力を出しやすくなる。
代償：リセットが急ぐ分だけ1次巻線の跳ね返り電圧が大きくなり，
\rev{スイッチに要求される耐圧が$2V_{\mathrm{in}}$から$3V_{\mathrm{in}}$へ上がる。}

%%--- 章末問題 6.4 ---
\item[\ref{ex6.4}]
\rev{周期は$T = 1/f = 50\,\mathrm{\mu s}$，オン時間は$DT = 15\,\mathrm{\mu s}$である。}
(a) \rev{式\ref{eq6.4}より，励磁電流はオン期間に傾き$V_{\mathrm{in}}/L_m$で増えるから}
\begin{equation*}
\rev{I_{\mathrm{pk}} = \frac{V_{\mathrm{in}}}{L_m}DT
= \frac{100}{10\times10^{-3}}\times15\times10^{-6} = 0.15\,\mathrm{A}}
\end{equation*}
(b)
\begin{equation*}
\rev{W = \frac{1}{2}L_m I_{\mathrm{pk}}^2
= \frac{1}{2}\times10\times10^{-3}\times0.15^2 \approx 0.11\,\mathrm{mJ}}
\end{equation*}
(c) 式\ref{eq6.14}より$t_r = (N_3/N_1)DT = 15\,\mathrm{\mu s}$。
オフ期間は$(1-D)T = 35\,\mathrm{\mu s}$だから，リセットは余裕をもって完了する
（式\ref{eq6.15}の条件$D \le 0.5$も満たしている）。
(d) 励磁電流を戻す道がなくなるので，毎周期$0.15$\,Aずつ積み上がる。
磁束が増え続け，鉄心は数周期のうちに磁気飽和して
励磁インダクタンスを失い，1次電流が急増してスイッチを破壊する。
トランスに「蓄えっぱなし」が許されないことの数値例である。

%%--- 章末問題 6.5 ---
\item[\ref{ex6.5}]
(a) \rev{式\ref{eq6.20}に$N_2/N_1 = 1$，$V_{\mathrm{out}} = V_{\mathrm{in}} = 12$\,Vを代入すると}
$\dfrac{D}{1-D} = 1$，よって$D = 0.5$。
(b) $N_2/N_1 = 2$より
\begin{equation*}
\rev{\frac{D}{1-D} = \frac{N_1}{N_2}\cdot\frac{V_{\mathrm{out}}}{V_{\mathrm{in}}}
= \frac{1}{2}\times\frac{48}{12} = 2}
\qquad\Rightarrow\qquad D = \frac{2}{3} \approx 0.67
\end{equation*}
昇圧の一部を巻数比（$\times2$）が受け持つことで，
$D$を極端に大きくせずに4倍の昇圧ができている。

%%--- 章末問題 6.6 ---
\item[\ref{ex6.6}]
(a) 式\ref{eq6.20}を巻数比について解くと
\begin{equation*}
\frac{N_1}{N_2} = \frac{D}{1-D}\cdot\frac{V_{\mathrm{in}}}{V_{\mathrm{out}}}
= \frac{0.4}{0.6}\times\frac{141}{5} \approx 18.8
\end{equation*}
(b) 巻数は整数だから$N_1 : N_2 = 19 : 1$に丸める。式\ref{eq6.20}より
\begin{equation*}
\frac{D}{1-D} = \frac{N_1}{N_2}\cdot\frac{V_{\mathrm{out}}}{V_{\mathrm{in}}}
= 19\times\frac{5}{141} \approx 0.674
\qquad\Rightarrow\qquad
D \approx 0.40
\end{equation*}
となり，ほぼ狙いどおりの$D$で動作できる。
$141$\,Vから$5$\,Vという28倍の変換でも，
巻数比が大部分を受け持つおかげで$D$は0.4に保てる。
ACアダプタや携帯機器の充電器で
フライバックが定番である理由がここにある。

%%--- 章末問題 6.7 ---
\item[\ref{ex6.7}]
(a) 理論値は式\ref{eq6.20}より
\begin{equation*}
\rev{V_{\mathrm{out}} = \frac{0.2941}{1-0.2941}\times12 \approx 5.0\,\mathrm{V}}
\end{equation*}
シミュレーションでも平均約$5$\,V（小さなリプル付き）となり一致する。
(b) \rev{$D = 0.5$では$V_{\mathrm{out}} \approx 12$\,V（等倍），
$D = 0.7059$，$V_{\mathrm{in}} = 5$\,Vでは$V_{\mathrm{out}} \approx 12$\,V（2.4倍の昇圧）となる。}
同じ回路がデューティ比だけで降圧にも昇圧にもなることが確認できる。
(c) スイッチが切れる瞬間，1次電流はピーク値から一気にゼロへ落ち，
\rev{同時に2次電流が$({N_1}/{N_2})I_{\mathrm{pk}}$（$L_1 = L_2$なら同じ値）から}
流れ始める。電流の担い手が1次巻線から2次巻線へ乗り換えても，
磁束（巻数$\times$電流）は連続であることが波形から読み取れる。

\end{enumerate}
